% =====================================================================
% Section 4 --- One kernel, five closures, three financings
% ---------------------------------------------------------------------
% Job: the common kernel (E in passing: one production core deliberately
%      held fixed); row and cell definitions; the "why only three economies"
%      box (Propositions 1-3 in one paragraph, forward reference to Sec. 6
%      and App. A); the presentation rule as one paragraph.
% Mindmap: N02A, N06, N11, N12.
% Status: STUB --- narrative outline v3 (2026-09-20).
% =====================================================================
\section{One kernel, five closures, three financings}\label{sec:kernel}

% TODO: verbal model overview BEFORE the equations (R1.2); state the static
% horizon explicitly; one production core (theory of production E) held fixed
% by design.

We examine how different closures have different effects.
For this we choose five different labor closures and three financing closures that reflect different assumptions about the economy.
The labour closures specify how wages and employment adjust.
\textbf{BF} fixes employment in each sector at its baseline level and lets
sector-specific wages adjust. \textbf{ALPHA} allows labour to move between
sectors at one flexible wage while keeping total employment fixed.
\textbf{BETA} lets sectoral employment respond to each sector's real wage,
deflated by the consumption price index, with a common supply elasticity
but separate wages. \textbf{GAMMA} fixes the common real wage and lets
employment adjust without a cap. \textbf{DELTA} applies GAMMA's wage rule
in the Leontief limit, with fixed proportions in production and consumption;
it is a corner rather than a separate labour-market mechanism.

The financing closures specify how programme-related demand is accommodated
and who bears its cost. \textbf{F1} redirects household spending towards
programme sectors by changing preferences: larger spending shares there are
paid for by smaller shares elsewhere within the household budget, without
additional public purchases. \textbf{F2} adds public investment financed by
a household tax equal to the programme's current-price cost. \textbf{F3}
adds the same public-investment bundle with external financing and no
additional domestic programme tax. This external financing is an accounting
counterpart, not a monetary mechanism.

\subsection{The five labour closures}\label{subsec:labour-closures}
% Formulations: registry/closures.toml; ADR-0020 and ADR-0022.
% Appendix A resets the equation counter; keep these PDF anchors distinct.

\begingroup
\renewcommand*{\theHequation}{labour.\arabic{equation}}
Let $L_i$ denote employment in sector $i$, $\bar{L}_i$ its baseline value,
and $\Lbar = \sum_{i=1}^N \bar{L}_i$ total baseline employment. Sectoral
wages are $w_i$; with mobile labour they equal one common wage $w$.
The real wage is deflated by the consumption price index
\begin{equation}
  \Pi(p) = \left(\sum_{i=1}^N \omega_i p_i^{1-\sigma}\right)^{1/(1-\sigma)},
  \qquad \sum_{i=1}^N \omega_i = 1,
  \label{eq:labour-cpi}
\end{equation}
with calibrated weights $\omega_i$ and the Cobb--Douglas limit at
$\sigma = 1$. Subscript $0$ denotes the baseline. All closures use the
same cost-minimizing labour demand,
\begin{equation}
L_i^{\mathrm{cm}}(p,y_i,w_i)
 = \alpha_i A_i^{\epsilon-1} y_i
   \left(\frac{p_i}{w_i}\right)^{\epsilon}.
\label{eq:closure-labour-demand}
\end{equation}
Here $y_i$ is gross output, $A_i$ productivity and $\alpha_i$ the calibrated
labour share. The elasticity $\epsilon$ governs substitution between labour
and the intermediate composite. The closures differ in the wage and
labour-supply restrictions imposed on this demand.

\paragraph{BF: immobile labour.}
The BF row fixes employment in every sector and solves a sector-specific
wage to clear each labour market:
\begin{equation}
L_i = L_i^{\mathrm{cm}}(p,y_i,w_i) = \bar{L}_i,
\qquad i=1,\ldots,N, \qquad \eta = 0.
\label{eq:closure-bf}
\end{equation}
Thus $\sum_i L_i = \Lbar$, but wages may differ across sectors.
The allocation parameter $\eta$ selects only the endpoints $\{0,1\}$;
its mobile endpoint $\eta = 1$ is ALPHA below, not an additional closure.

\paragraph{ALPHA: mobile labour at full employment.}
Labour reallocates across sectors at one flexible wage, while total
employment remains fixed:
\begin{equation}
L_i = L_i^{\mathrm{cm}}(p,y_i,w),
\qquad \sum_{i=1}^N L_i = \Lbar,
\qquad \eta = 1.
\label{eq:closure-alpha}
\end{equation}

\paragraph{BETA: elastic sectoral labour supply.}
We use the uniform-elasticity sectoral form: each sector has its own wage
and supply curve, with one common labour-supply elasticity $\eta_s$:
\begin{equation}
L_i = L_i^{\mathrm{cm}}(p,y_i,w_i)
 = \bar{L}_i \left[\frac{w_i/\Pi(p)}{w_{i0}/\Pi_0}\right]^{\eta_s},
\qquad \eta_s \geq 0, \qquad i=1,\ldots,N.
\label{eq:closure-beta}
\end{equation}
Equal supply elasticities do not impose equal wages: $N$ separate labour
markets determine $w_i$. Here $\eta_s$ is a labour-supply elasticity, not
the allocation parameter $\eta$; setting $\eta_s = 0$ gives exactly the
immobile BF row. The calibrated real-wage anchors are $w_{i0}/\Pi_0 = 1$.
The deflator remains the common consumption price index, not the sector's
own goods price $p_i$. The supply curves are reduced forms without a
labour--leisure income effect. The headline panels use $\eta_s = 0.5$;
the elasticity ladder and heterogeneous sectoral variants are reported in
Appendix~\ref{app:variants}. The single-wage BETA is retained there only
as a distinct benchmark; its zero-elasticity limit is ALPHA, not BF.

\paragraph{GAMMA: fixed real wage.}
The real wage is pegged at its baseline value, while employment is
endogenous and uncapped:
\begin{equation}
\frac{w}{\Pi(p)} = \bar{w} = 1,
\qquad L_i = L_i^{\mathrm{cm}}(p,y_i,w),
\qquad \eta = 1.
\label{eq:closure-gamma}
\end{equation}
Unlike ALPHA, this closure does not impose $\sum_i L_i = \Lbar$.
It is a two-sided wage peg, not a one-sided wage floor with unemployment
rationing. On the demand-only baseline branch, $\Pi(p)=1$ and the
implemented wage pin is $w=1$.

\paragraph{DELTA: the Leontief corner.}
DELTA retains GAMMA's wage and employment rule while taking the
substitution elasticities to their Leontief limits:
\begin{equation}
\frac{w}{\Pi(p)} = \bar{w} = 1,
\qquad \eta = 1,
\qquad (\theta,\epsilon,\sigma) \longrightarrow (0^+,0^+,0^+).
\label{eq:closure-delta}
\end{equation}
Here $\theta$ governs substitution between intermediate inputs and
$\sigma$ substitution in consumption. DELTA is therefore a corner of
GAMMA with fixed-coefficient technology, rather than an independent
labour-market mechanism.
\endgroup

\subsection{The three financing closures}\label{subsec:financing-closures}
% Formulations: registry/closures.toml; ADR-0002, ADR-0019 and ADR-0020.
% Programme-scale F1 tilt: experiments/run.jl, f1_tilt_weights.
\begingroup
\renewcommand*{\theHequation}{financing.\arabic{equation}}
Each labour closure is crossed with three ways of accommodating the same
programme scale and sectoral incidence. Let $g_i = kG_0\psi_i$ denote the
reference programme bundle, where $G_0$ is its baseline-price value,
$k \geq 0$ scales the programme, and $\psi_i \geq 0$ with
$\sum_i\psi_i = 1$ describes its sectoral incidence. F2 and F3 add this
bundle to final demand; F1 instead uses its scale and incidence to redirect
household spending.

Baseline government purchases $g_i^G$ remain tax-financed in all three
cases. Define the wage bill and their current-price tax cost as
\begin{equation}
W = \sum_{i=1}^N w_i L_i,
\qquad T_G(p) = \sum_{i=1}^N p_i g_i^G.
\label{eq:financing-baseline-budget}
\end{equation}
With a common wage, $W = w\sum_i L_i$. Let $E$ denote after-tax household
resources, including the net external transfer $F$, and let $s$ be the
calibrated saving rate. Household consumption expenditure is
$E^h = (1-s)E$. Gross household purchases $c_i^h$ exhaust this budget:
\begin{equation}
\sum_{i=1}^N p_i c_i^h = E^h,
\qquad c_i^{\mathrm{dom}} = (1-m_i)c_i^h,
\label{eq:financing-household-budget}
\end{equation}
where $m_i$ is the sectoral import margin. Only the domestic content enters
domestic goods-market clearing; likewise, the programme contributes
$(1-m_i)g_i$ under F2 and F3.

\paragraph{F1: budget-neutral preference reallocation.}
F1 changes the household's CES weights without adding a public-demand
bundle. For positive preference shifters $d_i$, the weights and gross
household demand become
\begin{align}
\widetilde{\omega}_i
 &= \frac{\omega_i d_i}{\sum_j\omega_j d_j},
\qquad \sum_i\widetilde{\omega}_i = 1,
\label{eq:closure-f1-weights}\\
c_i^h
 &= E^h\frac{\widetilde{\omega}_i p_i^{-\sigma}}
 {\sum_j\widetilde{\omega}_j p_j^{1-\sigma}},
\qquad E = W-T_G(p)+F.
\label{eq:closure-f1}
\end{align}
The programme-scale tilt is $d_i = 1+kG_0\psi_i^+/c_{i0}^h$ in sectors
with positive baseline household purchases, and $d_i=1$ otherwise;
$\psi^+$ is the programme incidence restricted to those sectors and
renormalized to sum to one. The CES normalizer preserves
\eqref{eq:financing-household-budget} at every price vector. Thus increased
spending shares in programme sectors are paid for by reduced shares
elsewhere, conditional on the household budget; equilibrium income need
not remain unchanged. F1 is a demand-composition experiment, not an
autonomous investment injection or multiplier.

\paragraph{F2: tax-financed public investment.}
F2 adds the real programme bundle $g$ and levies an additional household
tax equal to its value at equilibrium prices:
\begin{equation}
T(p) = \sum_{i=1}^N p_i g_i,
\qquad E = W-T_G(p)-T(p)+F,
\qquad B_{\mathrm{gov}}=0.
\label{eq:closure-f2}
\end{equation}
Total domestic tax revenue $T_G(p)+T(p)$ therefore balances the purchases
$g^G+g$, priced at current rather than baseline prices. The tax is
implemented as the endogenous wage-income rate
$\tau=[T_G(p)+T(p)]/W$; with one representative household its burden is
effectively lump-sum. The programme has an explicit domestic budget
counterparty and reduces household resources directly by $T(p)$.

\paragraph{F3: externally financed public investment.}
F3 adds the same real bundle but imposes no additional domestic programme
tax. Its current-price cost is booked as external programme financing:
\begin{equation}
B_{\mathrm{gov}} = \sum_{i=1}^N p_i g_i,
\qquad E = W-T_G(p)+F.
\label{eq:closure-f3}
\end{equation}
The booking $B_{\mathrm{gov}}$ is distinct from the net external transfer
$F$ entering household resources after tax. The latter is solved
endogenously in BF, ALPHA and both BETA forms, but is fixed at zero in
GAMMA and DELTA. With no unfinanced demand injection, all closures clear
all $N$ goods markets and satisfy the external-account identity
\begin{equation}
S + T_{\mathrm{int}} + M - (I+X) = F+B_{\mathrm{gov}},
\qquad S=sE,
\label{eq:financing-external-account}
\end{equation}
where $I$ is non-programme investment, $X$ exports, $M$ total imports
(including intermediate and programme imports), and $T_{\mathrm{int}}$
the separately booked product-tax leak on intermediate use; all are
current-price values. Here $B_{\mathrm{gov}}=0$ under F1 and F2, so the net
external position is $F+B_{\mathrm{gov}}$, not the programme cost alone.
F3 is an accounting open-economy closure, not a monetary mechanism: the
static model cannot distinguish external borrowing from a transfer and
does not model interest payments, repayment or exchange-rate adjustment.
\endgroup

\subsection{Why only three economies}\label{subsec:why-three}
% TODO: the equivalence box --- Propositions 1-3 in one paragraph
% (demand-only scope, the A = 1 branch the numeraire selects); forward
% reference to Sec. 6 and App. A.

\subsection{Presentation rule}\label{subsec:presentation-rule}
% TODO: ONE paragraph. Columns F1, F2, F3; shade F2 = F3 in the BF and
% ALPHA/BETA rows only; state the neutrality condition (F endogenous); flag
% the fixed-wage row as the financing-sensitive case; restate the collapse as
% fifteen cells -> seven economies (NOT three; three is a count of labour
% classes). Reference paper/framing_gamma.tex. Keep this as one paragraph ---
% do not leak it as a debate.

% ===== CARRIED OVER FROM AN OLD DRAFT =============================
% [devised for an old draft; not in the validated submission]
% source: revised_manuscript/chapters/ch03_cge_model.tex (commit 41144dd)
% =================================================================
% Eventually, the question emerges how to link such a setup with the Input-Output tables introduced in the beginning of section \ref{sec:IO_models}. In practice, such a link necessitates that $\Omega$ plays a role in the exact specification of the production function and/or consumption function. For a Cobb-Douglas function, as the most widely used specification, the cost shares contained in $\Omega$ are interpreted as output elasticities and, correspondingly, used to specify the relative contribution of factors and intermediates to net output of a given sector, which translates into
% and choose the CES consumer price index as the numeraire:$$
% \text{CPI} = \left( \sum_{u=1}^n \omega_{0u} p_u^{\theta  -1} \right)^{\frac{1}{\theta - 1}}.
% $$
% Calculating GDP directly using these post-shock prices would yield a nominal GDP measure. To obtain real GDP, one typically employs a price index such as the Laspeyres or Paasche index. We use the Laspeyres index, so that GDP is computed as
% $$
% \text{GDP} = \frac{\sum_{s=1}^n p_s \post c_s}{\sum_{s=1}^n p_s  c_s}
% $$
% which measures the change in consumption valued at pre-shock prices, providing a meaningful comparison of real economic activity before and after the shock.
% While this basic formal setup obviously differs from the Leontief model presented in Section~\ref{sec:leontief} in many details,
% a major practical difference in applying these models also emerges from the underlying paradigmatic perspectives:
% loosely spoken, the Leontief model follows a 'demand-side view', where production structures are assumed constant and final demands change exogenously. In contrast, general equilibrium approaches will start from a 'supply-side view', where economic changes arise from changes in the exogenous parameters -- like production functions, technology or endowments -- and demand will respond.
% NOTE: the closing 'demand-side vs supply-side view' sentence was rejected by Reviewer 2 (R2.7) as nonsensical in GE; keep only if reworked.
